% a04_Bh_bulk4  Proof of app:a04:lem:B2, part 4: the averaging vertices, structure and line lemmas, the sup bound

\subsubsection*{B2.7 The averaging vertices: structure and the photon slices}

$\mathfrak a_x$ is the order-$e^2$ contribution to $\langle r_x^2\rangle^{\fl}$ of the $A$-dependence of King's block averaging, the vertices \cite[(3.24)--(3.26)]{K1} with the contour convention \cite[(2.19), (3.27)]{K1}. We work at unit block scale ($t=5^{-j}$, $\Lambda_t$, $m',\mu'\in(0,m_0]$, $b=\mathfrak a_j$) and write $\kappa_0\defeq\alpha_0/\sqrt3$, so that \cref{app:a04:lem:CT} at unit block scale reads $|G_\tau(x,y)|\le C\tau^{-3}e^{-\kappa_0|x-y|}$ and, at block side $\ell$, $|D_i(u,v)|\le C\ell^2\tau^{-3}e^{-\kappa_0|u-v|/\ell}$.

\emph{Photon slices.} For $0\le i\le j$ let $\ell_i\defeq5^it$, $P_i$ the uniform block average onto the blocks of side $\ell_i$ (corner blocks \cite[(2.12)]{K1}), $D_i\defeq(-\Delta_t+\mu'^2+\mathfrak a_i\ell_i^{-2}P_i)^{-1}$ ($D_0\defeq0$), and $S_i\defeq D_{i+1}-D_i$ for $i<j$. By \cite[(2.17)--(2.18)]{K1} the photon fluctuation field after $j$ steps is $A'=\sum_{i<j}A'_i$ with independent Gaussian slices, $\langle A'_{i,\mu}(u)A'_{i',\nu}(v)\rangle=\delta_{ii'}\delta_{\mu\nu}S_i(u,v)$ (\cref{app:a04:def:bulk}: $D_j=\sum_{i<j}G^\ph_{(i)}$).

\emph{Contours.} By \cite[(2.19)]{K1}, $A(\Gamma_{y,x})=\sum_{i<j}A'_i(\Gamma^{(i+1)}_{x_{i+1},x})$, $A(\Gamma)\defeq\sum_{\mathsf b\in\Gamma}tA_{\mathsf b}$, where $x_{i+1}$ is the corner of the $\ell_{i+1}$-block of $x$ and $\Gamma^{(i+1)}_{x_{i+1},x}$ is a staircase of $\le3\sum_{i'\le i}5^{i'+1}$ fine bonds, hence of length $\le C_L\ell_i$, $C_L\defeq3L^2/(L-1)$: slice $i$ sees only the sub-contour inside its own $\ell_{i+1}$-block.

\emph{Vertices.} With $U(\vartheta)=1+\vartheta J-\vartheta^2/2+O(\vartheta^3)$ ($J$ the complex structure on the two real components) and $V_{\mathrm{av}}=-\frac b2\sum_y[|Q(A)\phi|^2-|Q\phi|^2]$, $(Q(A)\phi)(y)=\sum_{x\in B(y)}t^3U(e'A(\Gamma_{y,x}))\phi(x)$ \cite[(3.24)--(3.26)]{K1} at zero block fields:
\begin{gather*}
V_1=-be'\sum_y\sum_{x\in B(y)}t^3A(\Gamma_{y,x})\,\ip{Q\phi(y)}{J\phi(x)},\qquad
V_2=\tfrac{be'^2}2\sum_{y,x}t^3A(\Gamma_{y,x})^2\ip{Q\phi(y)}{\phi(x)}-\tfrac{be'^2}2\sum_y\Big|\sum_{x\in B(y)}t^3A(\Gamma_{y,x})J\phi(x)\Big|^2,\\
V_h=e't^3\sum_{\mathsf b}A_{\mathsf b}\ip{\nabla\phi(\mathsf b)}{J\phi(\mathsf b_+)}=-e't^2\sum_{\mathsf b}A_{\mathsf b}\ip{\phi(\mathsf b_-)}{J\phi(\mathsf b_+)},
\end{gather*}
the last being the hopping vertex \cite[(3.22)]{K1} at $m=1$ (one difference quotient) at unit block scale ($\ip{\phi(\mathsf b_+)}{J\phi(\mathsf b_+)}=0$; its sign plays no role). At order $e'^2$ of $\langle|\phi(x_0)|^2\rangle^{\fl}$, since $\langle A'\rangle=0$ kills the first order and everything without an averaging vertex is the original-vertex part, $\mathfrak a_{x_0}$ consists of the four classes of connected correlations
\[
\mathrm{E1}+\mathrm{E2}=\langle|\phi(x_0)|^2;V_2\rangle,\qquad\mathrm{E3}=\tfrac12\langle|\phi(x_0)|^2;V_1;V_1\rangle,\qquad\mathrm{E4}=\langle|\phi(x_0)|^2;V_1;V_h\rangle.
\]
Their photon kernels are $\cV(x)\defeq\sum_i\Var(A'_i(\Gamma^{(i+1)}_{x_{i+1},x}))$, $\cK(x,x')\defeq\sum_i\langle A'_i(\Gamma(x))A'_i(\Gamma(x'))\rangle$ and $\cK_{\mathsf b}(x)\defeq\sum_i\cK_{i,\mathsf b}(x)$, $\cK_{i,\mathsf b}(x)\defeq\sum_{\mathsf b'\in\Gamma_i(x),\ \mathsf b'\parallel\mathsf b}tS_i(\mathsf b',\mathsf b)$; the scalar lines are $G=G_t$, $(QG)$, $(QGQ^*)$. Each class is a two- or three-line cycle through $x_0$ (E1/E2: two lines from $x_0$ to legs in one block; E3/E4: $x_0$--$v_1$--$v_2$--$x_0$), and E4 carries one differentiated line: the $J$-antisymmetry at $\mathsf b$ gives $2(QG)(y,x_0)[G(x_0,\mathsf b_-)G(x,\mathsf b_+)-G(x_0,\mathsf b_+)G(x,\mathsf b_-)]$, a forward difference at $\mathsf b$ on one of the two lines.

\begin{lemma}[photon slices]\label{app:a04:lem:S}
\begin{enumerate}[label=\textup{(S\arabic*)},leftmargin=*,start=0]
\item $S_i$ is the slice covariance of \cite[(2.17)]{K1}: the joint Gaussian of $(\phi,\psi_i)$ there has marginal $\phi$-covariance $D_{i+1}$, because $\ell_{i+1}^2/\mathfrak a_{i+1}=\ell_i^2/\mathfrak a_i+\ell_{i+1}^2/\mathfrak a$ (algebra from \cite[(2.13)]{K1}), and by the law of total covariance $D_{i+1}=\Cov(\phi|\psi)+\Cov(\E[\phi|\psi])=D_i+\mathfrak a_i^2\ell_i^{-4}D_iQ_i^*C^{(i)}Q_iD_i$, the term of \cite[(2.17)]{K1}.
\item $0\le S_i\le D-D_i$: $P_{i+1}=Q_i^*Q'^*Q'Q_i\le P_i$ ($Q'Q'^*=1$) and $\mathfrak a_i\ell_i^{-2}\ge\mathfrak a_{i+1}\ell_{i+1}^{-2}$, so $D_i\le D_{i+1}\le D$.
\item $|S_i(u,v)|\le C\ell_i^{-1}$: $S_i(u,u)\le(D-D_i)(u,u)$ by (S1), Cauchy--Schwarz, and $(D-D_i)(u,u)=\ell_i^{-1}\Gamma_{t/\ell_i}(u/\ell_i,u/\ell_i)\le C\ell_i^{-1}$ by \cref{app:a04:lem:B0}(i) at relative spacing $t/\ell_i=5^{-i}$, mass $\mu'\ell_i\le m_0$, coupling $\mathfrak a_i$; for $i=0$, $D(u,u)\le C/t$. Uniform in $\mu'$.
\item $|\nabla_uS_i(u,v)|\le C\ell_i^{-2}$: ($\Gamma$2) rescaled for $D-D_i$ and $D-D_{i+1}$ ($i\ge1$), (T1) for $i=0$.
\item $|S_i(u,v)|\le C\ell_i^{-1}\min\big(1,e^{(\kappa_0/2)(\rho_i-|u-v|)/\ell_{i+1}}\big)$ with $\rho_i\defeq(6/\kappa_0)\ell_{i+1}\ln(\ell_{i+1}/t)\le C\ell_i(i+1)$: \cref{app:a04:lem:CT} for $D_i$ and $D_{i+1}$ at block side $\le\ell_{i+1}$, half the exponent paying the prefactor $\ell_{i+1}^2t^{-3}$ beyond $\rho_i$.
\item For $u,v\in\Lambda_t\subset\Lambda_{t/5}$ and $i\ge1$: $|S^{(t)}_i(u,v)-S^{(t/5)}_{i+1}(u,v)|\le Ct\ell_i^{-2}$: the coarse slice $i$ and the fine slice $i+1$ have the same block sides $\ell_i,\ell_{i+1}$; compare each of the four coarse lines with the continuum one at its block side by \cref{app:a04:lem:B0}(ii) rescaled ($\ell^{-1}\cdot C(t/\ell)$), and the coupling change $\mathfrak a_{i}\to\mathfrak a_{i+1}$ by \cref{app:a04:lem:B0}(iv) rescaled ($C\ell_i^{-1}5^{-2i}=Ct^2\ell_i^{-3}\le Ct\ell_i^{-2}$). Slice $0$ of both members and the extra finest slice of the finer member are not matched; each contributes $O(t\cdot\mathrm{polylog})$ (E4: $O(t^2\cdot\mathrm{polylog})$) by \cref{app:a04:lem:C}(i),(ii) below, since $\ell_0=t$ for both.
\end{enumerate}
\end{lemma}

\begin{lemma}[contour kernels]\label{app:a04:lem:C}
\begin{enumerate}[label=\textup{(\roman*)},leftmargin=*]
\item $\Var(A'_i(\Gamma^{(i+1)}))\le(C_L\ell_i)^2C\ell_i^{-1}=C\ell_i$; $\cV\le C\sum_{i<j}\ell_i\le C$; $|\cK_i|\le C\ell_i$ and $|\cK|\le C$ (Cauchy--Schwarz).
\item $|\cK_{i,\mathsf b}(x)|\le C$ ($\le C_L\ell_i/t$ bonds of weight $t$, (S2)); up to the tail of (S4) it is supported in $|\mathsf b-x|\le r_i\defeq C_L\ell_i+\rho_i\le C\ell_i(i+1)$.
\item For $x\in\Lambda_t$ (corner blocks, $x_1=x$) the finer member's sub-contours coincide with the coarse ones as point sets and its finest slice sees the empty contour; for $x'$ in the coarse cell of $x$ every slice gains $\Gamma_{x,x'}$ ($\le3(L-1)$ fine bonds, length $\le3t$).
\item Replacing the fine bond sum along a contour by the coarse one in the same kernel costs $\le2(C_L\ell_i)^2C\ell_i^{-2}t=Ct$ per slice ((S3)).
\item In-cell oscillation of the finer member's kernels: $|\cK'(x',\cdot)-\cK'(x,\cdot)|\le Ct(j+1)$; for $\cK'_{\mathsf b}$ the E4 cost is $\le Ctj^2$.
\end{enumerate}
\end{lemma}
\begin{proof}
(i): the variance of $A'_i$ along a contour of $\le C_L\ell_i/t$ bonds of length $t$ is bounded by (length)$^2\sup|S_i|$, (S2); $\sum_{i<j}\ell_i=t(5^j-1)/4\le\frac14$. (ii): (S2) and (S4), the bonds of $\Gamma_i(x)$ lying within $C_L\ell_i$ of $x$. (iii) is the corner-block convention \cite[(2.12), (2.19)]{K1}. (iv): (S3), the Riemann error of a bond sum with step $t$ over a path of length $C_L\ell_i$ against a kernel with gradient $\le C\ell_i^{-2}$. (v): (iii) and (S2), (S3), summed over the $\le j+1$ slices.
\end{proof}

\begin{lemma}[scalar lines]\label{app:a04:lem:G}
Write $|r|_t\defeq|r|+t$ and $R_t\defeq(6/\kappa_0)\ln t^{-1}=6\sqrt3\,\alpha_0^{-1}\ln t^{-1}$, the radius used in the proof of \cref{app:a04:lem:B1} (in \S\S B2.7--B2.10 $R_t$ denotes this radius, not the $R_\tau$ of (O2) in \cref{app:a04:lem:TIOUT}); $G=C_t-\Gamma_t$.
\begin{enumerate}[label=\textup{(G\arabic*)},leftmargin=*]
\item $|C_t(r)|\le C|r|_t^{-1}$, $|\nabla C_t|\le C|r|_t^{-2}$, $|\nabla\nabla C_t|\le C|r|_t^{-3}$ (T1); $|C_t-C_{t/5}|(r)\le Ct^2|r|^{-3}$ and $|\nabla^{(t)}C_t-\nabla^{(t/5)}C_{t/5}|(r)\le Ct|r|^{-3}$ for $r\in\Lambda_t\setminus0$ (\cref{app:a04:lem:LP,app:a04:lem:LPgrad} at both spacings against $C_0$, $\partial C_0$).
\item $|\Gamma_t|,|\nabla\Gamma_t|\le C$; $|\Gamma_t-\Gamma_{t/5}|\le Ct+C5^{-2j}$ (($\Gamma$4) and ($\Gamma$5) for the change of $b$); $|\partial_b\Gamma_t|\le C$; $|\nabla\nabla\Gamma_\tau(u,v)|\le C\sum_{n\in N_\tau}|p|^2\sigma\hat C_\tau\le C\sum_{|n|\le C/\tau}\rho(n)\le C(1+\ln\tau^{-1})^3$ (uniform in $m'$).
\item $|G(x,y)|\le Ct^{-3}e^{-\kappa_0|x-y|}$ (\cref{app:a04:lem:CT}); hence $|G|\le C|r|_t^{-1}+C$ on $|r|\le R_t$ and $\le Ct^3$ beyond; $|(QG)|,|(QGQ^*)|\le C$ with the same tail; $t^3G(x,x)\le Ct^2$.
\item Gradient line: $|\nabla G|\le C|r|_t^{-2}+C$; $|\nabla_{\mathsf b}(QG)(y,\mathsf b)|\le C$ uniformly in $y$.
\end{enumerate}
\end{lemma}
\begin{proof}
(G1), (G2) restate \cref{app:a04:lem:TIOUT,app:a04:lem:Gam,app:a04:lem:LP,app:a04:lem:LPgrad}; the second-difference bound on $\Gamma_\tau$ follows from the Fourier representation with two factors $|d_\tau(p_\mu)|\le|p|$ and (E1)--(E3). (G3), (G4): (G1), (G2) and \cref{app:a04:lem:CT}; $(QG)(y,\cdot)$ is a block average of $G$.
\end{proof}

\begin{lemma}[tails of the contour kernels]\label{app:a04:lem:T}
Fix $i\ge1$, $r\defeq|\mathsf b-x|$, $\lambda_i\defeq2\ell_{i+1}/\kappa_0$, $M\defeq(2/\kappa_0)\ln t^{-1}$ and $\varrho_i\defeq r_i+M\ell_{i+1}$; for $x,\mathsf b$ coarse let $\Delta\cK_{i,\mathsf b}(x)\defeq\cK^{(t)}_{i,\mathsf b}(x)-\cK^{(t/5)}_{i+1,\mathsf b}(x)$ (coarse slice $i$ matched with fine slice $i+1$).
\begin{enumerate}[label=\textup{(T\arabic*)},leftmargin=*]
\item $|\Delta\cK_{i,\mathsf b}(x)|\le Ct/\ell_i$ for every $r$ ((S5) and \cref{app:a04:lem:C}(iv): $\le(C_L\ell_i/t)\cdot t\cdot Ct\ell_i^{-2}$).
\item $|\cK^{(t)}_{i,\mathsf b}(x)|\le Ce^{-(r-r_i)/\lambda_i}$ and $|\cK^{(t/5)}_{i+1,\mathsf b}(x)|\le CL^3e^{-(r-r_i)/\lambda_i}$ for $r>r_i$, hence $|\Delta\cK_{i,\mathsf b}(x)|\le CL^3e^{-(r-r_i)/\lambda_i}$.
\item (E4 cost, both Wick families.) Let $\kappa\ge0$ be radial and non-increasing with $\kappa(r)\ge|\Delta\cK_{i,\mathsf b}(x)|\mathbf 1_{r>r_i}$, e.g.\ $\kappa_i(r)\defeq\min(Ct/\ell_i,CL^3e^{-(r-r_i)_+/\lambda_i})$. Then the E4 term with the kernel $\Delta\cK_{i,\mathsf b}\mathbf 1_{r>r_i}$ is $\le Ct(1+j)^3R_t^3$ per slice, and the same holds for the kernel $|\cK^{(\tau)}_{i,\mathsf b}|\mathbf 1_{r>r_i}$ of a single member with $Ct$ replaced by $C\ell_{i+1}$.
\end{enumerate}
\end{lemma}
\begin{proof}
(T2): bound each member separately. Coarse member: the contour $\Gamma^{(i+1)}_{x_{i+1},x}$ ends at $x$ and has length $\le C_L\ell_i$, so each of its bonds $\mathsf b'$ has $|\mathsf b'-\mathsf b|\ge r-C_L\ell_i$, and (S4) gives $|S_i(\mathsf b',\mathsf b)|\le C\ell_i^{-1}e^{(\kappa_0/2)(\rho_i+C_L\ell_i-r)/\ell_{i+1}}=C\ell_i^{-1}e^{-(r-r_i)/\lambda_i}$; $\le C_L\ell_i/t$ bonds of weight $t$ give $\le Ce^{-(r-r_i)/\lambda_i}$. Finer member (slice $i+1$ of $\Lambda_{t/5}$): same scale $\ell_i$, same block side $\ell_{i+1}$, contour length $3(t/5)\sum_{i'\le i+1}5^{i'+1}\le C_L\ell_i$; its (S4) shift is $(6/\kappa_0)\ell_{i+1}\ln(\ell_{i+1}/(t/5))=\rho_i+(6/\kappa_0)\ell_{i+1}\ln5$, which multiplies the bound by $e^{3\ln5}=L^3$.
(T3): E4 contracts $|\phi(x_0)|^2$ against $V_1\ni\ip{Q\phi(y)}{J\phi(x)}$ and $V_h\ni\ip{\phi(\mathsf b_-)}{J\phi(\mathsf b_+)}$ (self-loops vanish, $\Tr J=0$; $x_0$ joined to both fields of one vertex leaves a self-loop on the other), so there are two Wick families, both carrying the kernel $\cK_{i,\mathsf b}(x)$: (i) $x_0$--$Q\phi(y)$, $x$--$\mathsf b$, and (ii) $x_0$--$\phi(x)$, $Q\phi(y)$--$\mathsf b$. \emph{Family (i).} By (G3)--(G4), uniformly in $\mathsf b$, for both placements of the difference (undifferentiated or gradient line from $\mathsf b$ to $x$), $\sum_xt^3(|G(x,\mathsf b_+)|+|\nabla G(x,\mathsf b)|)\kappa(|x-\mathsf b|)\le C\int(r^{-2}+r^{-1}+1)\kappa(r)d^3r+Ct\sup\kappa$ (the $O(1)$ points with $|x-\mathsf b|\le2t$ cost $\le t^3Ct^{-2}$ each). The outer factors: $\sum_{\mathsf b}t^3(|\nabla G(\mathsf b,x_0)|+|G(x_0,\mathsf b)|)\le CR_t^3$ (beyond $R_t$ a difference quotient is $\le2t^{-1}Ct^{-3}e^{-\kappa_0(|r|-t)}$, total $\le Ct^2R_t^2$) and $|(QG)(y,x_0)|\le C$. With $r_i\le C\ell_i(i+1)$, $M\ell_{i+1}=(2L\ln L/\kappa_0)\ell_ij$ and $\ell_i\le1$: $\varrho_i\le C\ell_i(1+j)$ and $\lambda_i\le2L/\kappa_0$. Shell $r_i<r\le\varrho_i$: $\kappa\le Ct/\ell_i$ by (T1), $\int_{r\le\varrho}(r^{-2}+r^{-1}+1)d^3r\le C\varrho(1+\varrho)^2$, cost $\le C(t/\ell_i)\ell_i(1+j)(1+j)^2=Ct(1+j)^3$. Region $r>\varrho_i$: by (T2), $\kappa\le CL^3e^{-(\kappa_0/2)M}e^{-(r-\varrho_i)/\lambda_i}=CL^3t\,e^{-(r-\varrho_i)/\lambda_i}$, and $\int_{r>\varrho}(r^{-2}+r^{-1}+1)e^{-(r-\varrho)/\lambda}d^3r\le C\lambda(1+\varrho+\lambda)^2\le C(1+j)^2$: cost $\le Ct(1+j)^2$. Near-diagonal: $Ct\sup\kappa\le Ct$. \emph{Family (ii)} (lines $G(x,x_0)$, $G(x_0,\mathsf b_\mp)$, $(QG)(y,\mathsf b_\pm)$, $y$ the block of $x$): the $J$-antisymmetry at $\mathsf b$ puts one forward difference at $\mathsf b$ on $G(x_0,\cdot)$ or on $(QG)(y,\cdot)$, and $|(QG)(y,\mathsf b)|,|\nabla_{\mathsf b}(QG)(y,\mathsf b)|\le C$ (G3)--(G4). The line from $x$ now runs to $x_0$. Split the $x$-sum: on $|x-x_0|\ge|x-\mathsf b|$, $|G(x,x_0)|\le C|x-\mathsf b|_t^{-1}+C$ ($|\cdot|_t^{-1}$ non-increasing); on $|x-x_0|<|x-\mathsf b|$, $\kappa(|x-\mathsf b|)\le\kappa(|x-x_0|)$. Hence, uniformly in $\mathsf b$ and $x_0$, $\sum_xt^3|G(x,x_0)|\kappa(|x-\mathsf b|)\le2\sup_z\sum_xt^3(C|x-z|_t^{-1}+C)\kappa(|x-z|)\le C\int(r^{-1}+1)\kappa(r)d^3r+Ct\sup\kappa$, which for $\kappa=\kappa_i$ is dominated by the three regions above ($r^{-1}+1\le r^{-2}+r^{-1}+1$; the first region's integral already runs over the whole ball $r\le\varrho_i$). The outer factor $\sum_{\mathsf b}t^3(|\nabla G(\mathsf b,x_0)|+|G(x_0,\mathsf b)|)\le CR_t^3$ is family (i)'s. Same per-slice cost $Ct(1+j)^3R_t^3$. For a single member, (T2) alone gives $\kappa\le Ce^{-(r-r_i)/\lambda_i}$ on $r>r_i$ and $\int_{r>r_i}(r^{-2}+r^{-1}+1)e^{-(r-r_i)/\lambda_i}d^3r\le C\lambda_i(1+r_i+\lambda_i)^2\le C\ell_{i+1}(1+j)^2$. The same holds for any radial non-increasing $\kappa$.
\end{proof}

\subsubsection*{B2.8 The sup bound for the averaging part}

\begin{proposition}[the averaging part, sup bound]\label{app:a04:prop:AVa}
For $f=1$, all $k$ and $1\le j\le k-J_0$, there are $q$ and $C$ (uniform in $m',\mu'\in(0,m_0]$, $b$, $k$, $j$) with $\sup_x|\mathfrak a_x|\le Ce^2(1+\ln t^{-1})^q$ at unit block scale, i.e.\ $\norm{\mathfrak a}_\infty\le Ce^2\eta^{-1}(1+j)^q$.
\end{proposition}
\begin{proof}
Bound each class by $|\text{photon kernel}|\times|\text{lines}|$, with the vertex sums restricted to $|\cdot-x_0|\le R_t$; the tails cost $\le Ct^3R_t^q$ by (G3). \emph{E1:} $2be'^2\sum_xt^3\cV(x)(QG)(y,x_0)G(x,x_0)\le C\sum_{|x-x_0|\le R_t}t^3(|x-x_0|_t^{-1}+1)\le C(1+R_t)^3$ ($\cV\le C$, \cref{app:a04:lem:C}(i), $|(QG)|\le C$). \emph{E2, E3:} $|\cK|\le C$ and $\sum t^3t^3(|x_0-u|_t^{-1}+1)(|u-v|_t^{-1}+1)(|v-x_0|_t^{-1}+1)$ over the $R_t$-balls $\le C(1+R_t)^6$ (local integrability in three dimensions; coincident points cost $Ct^2$). \emph{E4:} per slice $i$, with \cref{app:a04:lem:C}(ii) in its support form, $\sum_{\mathsf b}t^3|\nabla G(\mathsf b,x_0)|\sum_{|x-\mathsf b|\le r_i}t^3|G(x,\mathsf b_+)|\,|(QG)(y,x_0)|\le Cr_i^2(1+r_i)(1+R_t)^3$ (the ball sum of $|G(x,\mathsf b_+)|\le C|x-\mathsf b|_t^{-1}+C$ over $|x-\mathsf b|\le r_i$ is $\le C(r_i^2+r_i^3)$), the tail $|\mathsf b-x|>r_i$ costs $\le C\ell_{i+1}(1+j)^2R_t^3$ (\cref{app:a04:lem:T}(T3), single member), summable over $i$, and the other placements of the difference and the second Wick family are identical ($|\nabla_{\mathsf b}(QG)(y,\mathsf b)|\le C$). Summing over $i<j$, $\sum_ir_i^2(1+r_i)\le C(1+j)^3$ (dominated by $i\approx j$). Hence $\sup_x|\mathfrak a_x|\le Ce^2(1+\ln t^{-1})^q$ at unit block scale, with $e'^2=e^2\eta$ and the factor $\eta^{-1}$ of the density; $q$ is not pinned, and only its existence is used.
\end{proof}
